% part: intuitionistic-logic
% chapter: introduction
% section: natural-deduction

\documentclass[../../../include/open-logic-section]{subfiles}

\begin{document}

\olfileid{int}{int}{ntd}

\olsection{Deduksi Natural}

Deduksi natural tanpa aturan $\FalseCl$ iku sistem !!{derivation}
standar kanggo logika intuisionistik. Aturan-aturane diandharake
maneh ing kene kanthi alasan saka interpretasi BHK. Saben aturan
ngidini kita nyimpulake yen nalika premis-premise nduweni konstruksi,
kesimpulane uga nduweni konstruksi.

Amarga !!{derivation}s deduksi natural nduweni asumsi kang durung
dibebasake, !!a{derivation} kanggo $!A$ saka asumsi kang
!!{undischarged}~$\Gamma$ kudu dianggep fungsi kang ngowahi konstruksi
kanggo kabeh $!B \in \Gamma$ dadi konstruksi kanggo~$!A$. Yen ana
!!a{derivation} kanggo~$!A$ tanpa asumsi kang !!{undischarged}, ana
konstruksi kanggo~$!A$ miturut interpretasi BHK. Nanging, ing
andharan iki~$\Gamma$ bakal diilangi saka notasi yen ora dibutuhake.

Asumsi~$!A$ dhewe iku !!a{derivation} kanggo~$!A$ saka asumsi kang
!!{undischarged}~$!A$. Iki cocog karo interpretasi BHK: fungsi
identitas ing konstruksi ngowahi saben konstruksi kanggo~$!A$ dadi
konstruksi kanggo~$!A$.

\subsection{Konjungsi}

\begin{defish}
\AxiomC{$!A$}
\AxiomC{$!B$}
\RightLabel{\Intro{\land}}
\BinaryInfC{$!A \land !B$}
\DisplayProof
\hfill
\begin{tabular}{l}
\AxiomC{$!A \land !B$}
\RightLabel{\Elim{\land}}
\UnaryInfC{$!A$}
\DisplayProof
\\[3ex]
\AxiomC{$!A \land !B$}
\RightLabel{\Elim{\land}}
\UnaryInfC{$!B$}
\DisplayProof
\end{tabular}
\end{defish}

Upamane kita nduweni konstruksi $N_1$, $N_2$ kanggo $!A_1$ lan $!A_2$
miturut urutane. Mula kita uga nduweni konstruksi kanggo
$!A_1 \land !A_2$, yaiku pasangan $\tuple{N_1, N_2}$.

Konstruksi kanggo $!A_1 \land !A_2$ miturut interpretasi BHK yaiku
pasangan $\tuple{N_1, N_2}$. Dadi upamane kita nduweni pasangan kaya
mangkono. Kita uga nduweni konstruksi kanggo saben konjung:
$N_1$ iku konstruksi kanggo~$!A_1$ lan $N_2$ iku konstruksi kanggo $!A_2$.

\subsection{Kondisional}

\begin{defish}
  \AxiomC{$\Discharge{!A}{u}$}
  \DeduceC{$!B$}
  \DischargeRule{$\Intro{\lif}$}{u}
  \UnaryInfC{$!A \lif !B$}
  \DisplayProof
\hfill
  \AxiomC{$!A \lif !B$}
  \AxiomC{$!A$}
  \RightLabel{$\Elim{\lif}$}
  \BinaryInfC{$!B$}
  \DisplayProof
\end{defish}

Yen kita nduweni !!a{derivation} kanggo~$!B$ saka asumsi kang
!!{undischarged}~$!A$, ana fungsi~$f$ kang ngowahi konstruksi
kanggo~$!A$ dadi konstruksi kanggo~$!B$. Fungsi iku uga konstruksi
kanggo $!A \lif !B$. Dadi, yen premis $\Intro\lif$ nduweni konstruksi
kang gumantung marang konstruksi kanggo~$!A$, kesimpulan $!A
\lif !B$ nduweni konstruksi.

Ing sisih liyane, upamane ana konstruksi~$N$ kanggo $!A$ lan $f$
kanggo $!A \lif !B$. Konstruksi kanggo $!A \lif !B$ yaiku fungsi kang
ngowahi konstruksi kanggo $!A$ dadi konstruksi kanggo~$!B$. Mula
$f(N)$ iku konstruksi kanggo~$!B$: kesimpulan $\Elim\lif$ nduweni
konstruksi.

\subsection{Disjungsi}

\begin{defish}
\begin{tabular}{l}
\AxiomC{$!A$}
\RightLabel{\Intro{\lor}}
\UnaryInfC{$!A \lor !B$}
\DisplayProof
\\[3ex]
\AxiomC{$!B$}
\RightLabel{\Intro{\lor}}
\UnaryInfC{$!A \lor !B$}
\DisplayProof
\end{tabular}
\hfill
\AxiomC{$!A \lor !B$}
\AxiomC{$\Discharge{!A}{n}$}
\DeduceC{$!C$}
\AxiomC{$\Discharge{!B}{n}$}
\DeduceC{$!C$}
\DischargeRule{\Elim{\lor}}{n}
\TrinaryInfC{$!C$}
\DisplayProof
\end{defish}

Yen kita nduweni konstruksi~$N_i$ kanggo~$!A_i$, kita bisa ngowahi
dadi konstruksi~$\tuple{i, N_i}$ kanggo~$!A_1 \lor !A_2$. Ing sisih
liyane, upamane kita nduweni konstruksi kanggo~$!A_1 \lor !A_2$,
yaiku pasangan $\tuple{i, N_i}$ kang $N_i$ iku konstruksi kanggo~$!A_i$,
uga fungsi $f_1$, $f_2$ kang ngowahi konstruksi kanggo $!A_1$, $!A_2$
miturut urutane dadi konstruksi kanggo~$!C$. Mula $f_i(N_i)$ iku
konstruksi kanggo~$!C$, yaiku kesimpulan saka~$\Elim\lor$.

\subsection{Absurditas}

\begin{defish}
  \AxiomC{$\lfalse$}
  \RightLabel{\FalseInt}
  \UnaryInfC{$!A$}
  \DisplayProof
\end{defish}

Yen kita nduweni !!a{derivation} kanggo~$\lfalse$ saka asumsi kang
!!{undischarged}~$!B_1$, \dots, $!B_n$, ana fungsi~$f(M_1,
\dots, M_n)$ kang ngowahi konstruksi kanggo $!B_1$, \dots, $!B_n$
dadi konstruksi kanggo~$\lfalse$. Amarga $\lfalse$ ora nduweni
konstruksi, ora bisa ana konstruksi kanggo kabeh $!B_1$, \dots, $!B_n$
bebarengan. Mula $f$ uga nduweni sipat mangkene: \emph{yen} $M_1$,
\dots, $M_n$ iku konstruksi kanggo $!B_1$, \dots, $!B_n$ miturut
urutane, \emph{mula} $f(M_1, \dots, M_n)$ iku konstruksi kanggo~$!A$.

\subsection{Aturan kanggo $\lnot$}

Amarga $\lnot !A$ ditetepake minangka $!A \lif \lfalse$, sejatine
kita ora perlu aturan kanggo~$\lnot$. Nanging yen diwenehake,
aturane kaya mangkene:

\begin{defish}
\AxiomC{$\Discharge{!A}{n}$}
\noLine
\DeduceC{$\lfalse$}
\DischargeRule{\Intro{\lnot}}{n}
\UnaryInfC{$\lnot !A$}
\DisplayProof
\hfill
\AxiomC{$\lnot !A$}
\AxiomC{$!A$}
\RightLabel{\Elim{\lnot}}
\BinaryInfC{$\lfalse$}
\DisplayProof
\end{defish}


\subsection{Tuladha \usetoken{P}{derivation}}

\begin{enumerate}
\item $\Proves !A \lif (\lnot !A \lif \lfalse)$, yaiku $\Proves !A
  \lif ((!A \lif \lfalse) \lif \lfalse)$
  \begin{prooftree}
    \AxiomC{$\Discharge{!A}{2}$}
    \AxiomC{$\Discharge{!A \lif \lfalse}{1}$}
    \RightLabel{\Elim\lif}
    \BinaryInfC{$\lfalse$}
    \DischargeRule{\Intro\lif}{1}
    \UnaryInfC{$(!A \lif \lfalse)\lif \lfalse$}
    \DischargeRule{\Intro\lif}{2}
    \UnaryInfC{$!A \lif ((!A \lif \lfalse) \lif \lfalse)$}
  \end{prooftree}

\item $\Proves ((!A \land !B) \lif !C) \lif (!A \lif (!B \lif !C))$
  \begin{prooftree}
    \AxiomC{$\Discharge{(!A \land !B) \lif !C}{3}$}
    \AxiomC{$\Discharge{!A}{2}$}
    \AxiomC{$\Discharge{!B}{1}$}
    \RightLabel{\Intro\land}
    \BinaryInfC{$!A \land !B$}
    \RightLabel{\Elim\lif}
    \BinaryInfC{$!C$}
    \DischargeRule{\Intro\lif}{1}
    \UnaryInfC{$!B \lif !C$}
    \DischargeRule{\Intro\lif}{2}
    \UnaryInfC{$!A \lif (!B \lif !C)$}
    \DischargeRule{\Intro\lif}{3}
    \UnaryInfC{$((!A \land !B) \lif !C) \lif (!A \lif (!B \lif !C))$}
  \end{prooftree}

\item $\Proves \lnot(!A \land \lnot !A)$, yaiku $\Proves (!A \land (!A
  \lif \lfalse))\lif \lfalse$
  \begin{prooftree}
    \AxiomC{$\Discharge{!A \land (!A \lif \lfalse)}{1}$}
    \RightLabel{\Elim\land}
    \UnaryInfC{$!A \lif \lfalse$}
    \AxiomC{$\Discharge{!A \land (!A \lif \lfalse)}{1}$}
    \RightLabel{\Elim\land}
    \UnaryInfC{$!A$}
    \RightLabel{\Elim\lif}
    \BinaryInfC{$\lfalse$}
    \DischargeRule{\Intro\lif}{1}
    \UnaryInfC{$(!A \land (!A \lif \lfalse))\lif \lfalse$}
  \end{prooftree}

\item $\Proves \lnot\lnot(!A \lor \lnot !A)$, yaiku $\Proves ((!A \lor
  (!A \lif \lfalse))\lif \lfalse)\lif \lfalse$
  \begin{prooftree}\hskip-3em
    \AxiomC{$\Discharge{(!A \lor (!A \lif \lfalse))\lif \lfalse}{2}$}
    \AxiomC{$\Discharge{(!A \lor (!A \lif \lfalse))\lif \lfalse}{2}$}
    \AxiomC{$\Discharge{!A}{1}$}
    \RightLabel{\Intro\lor}
    \UnaryInfC{$!A \lor (!A \lif \lfalse)$}
    \RightLabel{\Elim\lif}
    \BinaryInfC{$\lfalse$}
    \DischargeRule{\Intro\lif}{1}
    \UnaryInfC{$!A \lif \lfalse$}
    \RightLabel{\Intro\lor}
    \UnaryInfC{$!A \lor (!A \lif \lfalse)$}
    \RightLabel{\Elim\lif}
    \insertBetweenHyps{\hskip -4em}
    \BinaryInfC{$\lfalse$}
    \DischargeRule{\Intro\lif}{2}
    \UnaryInfC{$((!A \lor (!A \lif \lfalse))\lif \lfalse)\lif \lfalse$}
  \end{prooftree}
\end{enumerate}

\begin{prop}
  Yen $\Gamma \Proves !A$ ing deduksi natural kanggo logika
  intuisionistik, $\Gamma \Proves !A$ uga ing logika klasik.
  Mligine, yen $!A$ iku teorema intuisionistik, uga dadi teorema klasik.
\end{prop}

\begin{proof}
  Saben aturan deduksi natural intuisionistik uga aturan ing deduksi
  natural klasik. Mula saben !!{derivation} ing logika intuisionistik
  uga !!a{derivation} ing logika klasik.
\end{proof}

\begin{prob}
  Wenehana !!{derivation}s ing logika intuisionistik kanggo !!{formula}s iki:
  \begin{enumerate}
    \item $(\lnot !A \lor !B) \lif (!A \lif !B)$
    \item $\lnot\lnot\lnot !A \lif \lnot !A$
    \item $\lnot\lnot (!A \land !B) \liff (\lnot\lnot !A\land \lnot \lnot !B)$
    \item $\lnot (!A \lor !B) \liff (\lnot !A \land \lnot !B)$
    \item $(\lnot !A \lor \lnot !B) \lif \lnot (!A \land !B)$
    \item $\lnot\lnot ( !A \land !B) \lif  (\lnot\lnot !A \lor
    \lnot\lnot !B)$
  \end{enumerate}
\end{prob}

\end{document}
